\documentclass[12pt]{article}
\usepackage{amsmath, amssymb, geometry}
\geometry{margin=1in}
\setlength{\parindent}{0pt}



\title{In-Class Exercise: Single-Variable Optimization and Logarithms\\[7pt]
\large ECON 320 - Introduction to Mathematical Economics}
\author{Muhebullah Karimzada}
\date{Fall 2025}
\begin{document}

\maketitle

\vspace{1em}

\textbf{1. Maximization of Profit Function}

A firm's total revenue and total cost functions are given by:
\[
R(Q) = 80Q - 2Q^2, \quad C(Q) = 20 + 10Q
\]

\textbf{Tasks:}
\begin{enumerate}
  \item Find the profit function \( \pi(Q) \).
  \item Determine the output level \( Q^* \) that maximizes profit.
  \item Compute the maximum profit.
  \item Verify that the second-order condition for a maximum holds.
\end{enumerate}

\vspace{1em}

\textbf{2. Utility Maximization Problem}

An individual's utility function is:
\[
U(x) = 4x - 0.5x^2
\]
where \( x \) is the number of hours spent on leisure.

\textbf{Tasks:}
\begin{enumerate}
  \item Determine the value of \( x \) that maximizes utility.
  \item Check whether it is a maximum using the second derivative test.
  \item Interpret the economic meaning of your result.
\end{enumerate}

\vspace{1em}
\newpage
\textbf{3. Logarithmic Utility and Diminishing Marginal Utility}

Consider the utility function:
\[
U(x) = \ln(x)
\]

\textbf{Tasks:}
\begin{enumerate}
  \item Compute \( U'(x) \) and \( U''(x) \).
  \item Show that marginal utility is positive but diminishing.
  \item If \( x = 10 \), find the elasticity of utility with respect to \( x \):
  \[
  E_{Ux} = \frac{dU}{dx} \cdot \frac{x}{U}
  \]
\end{enumerate}

\vspace{1em}

\textbf{4. Production Function and Marginal Product}

A production function is given by:
\[
Q = 10L^{0.5}
\]
where \( Q \) is output and \( L \) is labour.

\textbf{Tasks:}
\begin{enumerate}
  \item Find the marginal product of labour (MPL).
  \item At what level of \( L \) does diminishing marginal productivity begin?
  \item Suppose the wage rate \( w = 2 \) and output price \( p = 5 \). Determine the profit-maximizing \( L^* \).
\end{enumerate}

\vspace{1em}

\textbf{5. Maximization with Logarithmic Revenue}

A monopolist faces the demand function:
\[
P(Q) = 100 - 4Q
\]
and the total cost function:
\[
C(Q) = 20Q + 10
\]
Define revenue as \( R(Q) = P(Q) \cdot Q \). The firm decides to maximize:
\[
\Pi(Q) = \ln[R(Q)] - C(Q)
\]

\textbf{Tasks:}
\begin{enumerate}
  \item Express \( R(Q) \) and \( \Pi(Q) \) explicitly as functions of \( Q \).
  \item Find \( Q^* \) that maximizes \( \Pi(Q) \).
  \item Interpret why taking logs of revenue might change the firm's optimal decision compared to maximizing \( R(Q) - C(Q) \).
\end{enumerate}

\vspace{1em}

\textbf{6. Elasticity and Logarithmic Differentiation}

If \( y = 20x^{0.8} \), use logarithmic differentiation to find the elasticity of \( y \) with respect to \( x \).

\textit{Hint:} Take the natural logarithm of both sides and differentiate with respect to \( x \).

\vspace{1em}

\textbf{7. Application of Natural Logarithm Properties}

Simplify the following and interpret economically:
\[
\ln \left( \frac{x^3 y^2}{z^4} \right)
\]

If \( x \) represents capital, \( y \) labour, and \( z \) land, explain what happens to the output elasticity if \( x \) doubles while \( y \) and \( z \) remain constant.

\vspace{1em}



\end{document}
